\part{Coding Guidelines}\label{sec:CodingGuidelines}
These coding guidelines are a collection of coding standards and best practices. Obeying them should make your code better readable and less error prone. Some of them will even help to increase the program's overall performance. So perhaps you should see them not as optional, but more as obligatory rules.

But as myself, you may have issues to accept something as a \bdq{best practice} just because some guru told so. That's a good approach! First because what's a \bdq{best practice} today can be an \bdq{anti-pattern} tomorrow, and not just because some fashion trends have changed, but also because some new technology was introduced, or – happens all too often! – someone unmasked a false guru.

Another reason is that too often the explanations are missing why the \bdq{best practice} is superior to the alternatives. I have tried to give that explanations when I describe a best practice, but sometimes I just quote one of the gurus of our business and their explanations.

Coding standards are something different; a standard means to do the same thing always in the same way, and that is definitely an \textit{undoubted best practice}. Nevertheless there could be better ways to do that thing – in that case, change the standard\footnote{And I would appreciate if you update me accordingly so that I can modify this document.}.

But as always there may be good reasons to do it different from what is recommended or even requested by the guidelines\footnote{You remember the bullet points~\ref{lst:ZoP:SpecialCases} and \ref{lst:ZoP:Practicality} from the \ngref? If not, here are they again: \bdq{\ref{lst:ZoP:SpecialCases}.~Special cases aren't special enough to break the rules~…} and \bdq{\ref{lst:ZoP:Practicality}.~…~although practicality beats purity}.}. In such case a comment is required that describes that reason.\footnote{Omitting this comment is also a deviance from the rule, requiring a comment to explain it. Also known as the Catch~22 \autocite{WIKIPEDIA:Catch22}.}

One basic recommendation is that you should not write the same code over and over again. This is also known as the \bdq{DRY Principle} (\bdq{Don't Repeat Yourself} \autocite{WIKIPEDIA:DRY}) and we will discuss this later again.

And please keep in mind that not always the shortest, most compact source code is the best. Also avoid what is known as \bdq{Premature Optimization}\footnote{Donald E. Knuth\index{Knuth, Donald E.} made the following statement on optimisation: \bdq{We should forget about small efficiencies, say about 97\% of the time: premature optimization is the root of all evil.} \autocite{Knuth:PrematureOptimization}}. Modern optimising compilers and run-time optimisers do a very good job to create compact object code, so in most cases the programmer can concentrate fully on writing readable and comprehensible code. In this context I would like to remind you on the quotation from Martin Fowler's\index{Fowler, Martin} book that I put in front of this document, and again to the \ngref, here the points \bdq{\ref{lst:ZoP:SimpleVsComplex}.~Simple is better than complex} and \bdq{\ref{lst:ZoP:ComplexVsComplicated}.~Complex is better than complicated}.

But although optimisation still may have some limits, comments will have never any impact on the runtime performance of a program. So please refer to chapter \tqfullvref{sec:CommentsWhen} (if not done already) and see the recommendations on when to apply comments to your code.

%------------------------------------------------------------------------------

\chapter{Error Handling}\label{sec:ErrorHandling}
A proper error handling is crucial for any code, no matter in which type of program it lives – but the definition of \textit{proper} depends significantly on the type of the final product, as it was defined \tqfullvref{sec:TypesOfProducts}.

%==============================================================================
% End of file!!
%==============================================================================


\chapter{Error Prevention}\label{sec:ErrorPrevention}
In the \hyperref[sec:ErrorHandling]{previous chapter} we discussed how to deal with errors (failures) when they occur. Some errors are caused by the environment (missing files, dead networks, unavailable external systems), others are caused by inappropriate handling of the software (illegal input values, wrong configurations, missing credentials). These problems are mostly unavoidable, at least from the view as a programmer. Of course you can design the UI of your product that it guides the users as much as possible to prevent them from doing things wrong, but at the very end, they are still humans (hopefully~…).

There are three other types of errors that can affect a software system:
\begin{enumerate}[nosep]
	\item{Syntax Errors}
	
	\item{Implementation or Logical Errors}
	
	\item{Design Error}
\end{enumerate}

For a programming language like \Java that requires a compilation step before a program can be executed, errors can be generally sorted into two categories: into compile time errors and in errors that surface during execution (usually referred to as \bdq{runtime errors}). Syntax errors clearly belong to the first category, and usually, they very easy to detect and in most cases easy to fix. In fact, you have to fix them, otherwise your code will never reaches production level.

Basically, this whole book is discussing techniques that should help you avoiding errors. Those described here are explicitly targeted to this goal, while the others are more implicit. Unfortunately is it difficult to provide recipes for design errors, as they often surface only after a long time (and often the design was originally right at the time it was made, but other design decisions made it wrong later).

Consequently, we are focussing to implementation and logical errors for the remainder of this chapter.

%------------------------------------------------------------------------------

\section{Swap Logic Errors for Compiler Errors}\label{sec:SwapToCompilerErrors}
The usual weapon to fight implementation and logical errors are (unit\index{unit test}) tests\index{test}. How to write good tests is an art in its own and deserves another book, but some basics are discussed in chapter \tqfullvref{sec:WritingTests}.

But some of these errors can be prevented by swapping them to compiler errors. From a book about Java programming\index{Simmons, Robert Jr.}
\begin{quote}
“One fundamental principle of programming is that, generally, it is best to swap a logic error for a compiler error. Compiler errors tend to be found in seconds and are corrected just as fast. Syntax errors are a good example […]

Logic errors, on the other hand, are the bane of all programmers. They hide and hate to reveal themselves. Logic errors seem to have minds of their own, constantly evading detection and dodging your efforts to pin down their cause. They can easily take a thousand times more effort to solve than the worst compiler errors. Worst of all, many logic errors are not found at all and occur only intermittently in sensitive places, which causes your customer to scream for a fix. Logic errors often require you to throw thousands of man-hours at them, only to finally discover that they are minor typos.”

Robert Simmons Jr.: \textit{Hardcore Java} \autocite{Simmons:HARDCORE_JAVA}
\end{quote}

\subsubsection{Principles}
\begin{enumerate}[label=P\arabic*.]
\item{Utilise the capabilities of your development tools to detect bugs as soon as possible.

Java is a strongly typed language, and this makes the compiler a powerful tool in the regard, given that you design your code in a manner that it can detect errors.}

\item{Take warnings serious!

The \verb#javac# in its standard configuration emits already several warnings, build tools like Maven\autocite{APACHE_MAVEN}\index{Apache Maven}\index{Maven|see{Apache Maven}} or Gradle\autocite{GRADLE}\index{Gradle} can add to them, and finally you can configure additional warnings in your favourite IDE.}

\item{Avoid to outsmart the compiler!

Casting an object to another type is a sample for this, invoking methods through Reflection\index{reflection} is another.\footnote{Thanks to the creators of Java that you cannot invoke a method through a pointer, as it is non uncommon in C and C++~…}}
\end{enumerate}

A quite simple example for swapping a potential logic error for a compiler error is the recommendation to write comparisons always with the unchangeable value on the left side\footnote{This recommendation originates from C/C++ programming where \lstinline|if| conditions are checking values of type integer – and in C, nearly everything can be an integer, or at least be interpreted like one. Java forces that the type of the expression in the \lstinline|if| condition has to be \lstinline|boolean|, therefore this approach is less useful for Java code.}. So if you forget the second equal sign for a comparison on equal, the compiler will complain immediately:
\begin{lstlisting}
if( length() == len ) …
if( 5 == len ) …
\end{lstlisting}

Another sample for this is the recommendation to name fields with the “\verb#m_#” prefix (refer to chapter \tqfullvref{sec:AttributeNames}) instead of using the \lstinline|this.| prefix. The compiler will never complain if you omit that prefix when accessing a field, and it will not complain if you name a local variable in the same way as a field – but it will scream loudly about a non-existing reference if you forget the “\verb#m_#” prefix when accessing a field named along the rules defined here. Although you can configure both Eclipse and IntelliJ Idea to raise an error or a warning if you access a field without \lstinline|this.|, or if you shadow a field through a local variable with the same name. Nevertheless, I prefer the “\verb#m_#” prefix for fields because it is also easier to read.

Also defining and using value types and domain value types instead of simple types can prevent logic errors from remaining undetected: providing an article number to a method that expects a customer id will not work, when you defined the domain value types \lstinline|ArticleId| and \lstinline|CustomerId| and declared the respective method to take an argument of type \lstinline|CustomerId| instead of a simple \lstinline|String| or \lstinline|int|. Refer to the chapter \tqfullvref{sec:ValueTypes} for more details about this topic.

For the same reason you should always consider to use enums instead of defining magic numbers\index{magic number}.

The best thing about compiler errors is, that they usually show up early during development – and not first after deployment at customer side. Under this aspect you should take compiler warnings seriously, no matter whether emitted by \verb#javac# or your development environment. When CI/CD pipelines are used, additional code checking tools may have been integrated. Utilise them, and respond to their warnings, at least by thinking twice before adding the \lstinline|@SuppressWarnings|\autocite{ORACLE_DOC_SUPPRESSWARNINGS_ANNOTATION}\index{java.lang!SuppressWarnings} annotation.

See also chapter \tqvref{sec:CompilerWarningsAndErrors} about compiler warnings and errors.

In this regard the use of \lstinline|java.lang.Class::forName|\autocite{ORACLE_DOC_CLASS:forName}\index{java.lang!Class!forName()} and the use of Reflection\index{reflection} in general are another nasty source for bugs that might be difficult to identify.

Basically, \lstinline|Class.forName()|\index{java.lang!Class!forName()} takes a string as its argument that is the fully qualified name of a class. If there is at least one instance of \lstinline|Class|\index{java.lang!Class} on the current classpath, the method will return that object – otherwise it will throw a \lstinline|ClassNotFoundException|\autocite{ORACLE_DOC_CLASSNOTFOUNDEXCEPTION_CLASS}\index{java.lang!ClassNotFoundException}. This will not happen at compile time, but only during the execution of the program, depending on the execution path of your program – meaning the bug that you introduced by having a typo in the class or package name may remain undetected for a long time.

So if you want to access a class object that you expect to be there, always use “\lstinline|<ClassName>.class|”\index{java.lang!Class!class} where \verb#<ClassName># is the name of the class; you can use the fully qualified name, or you can import it. Either way, the compiler will complain if the class is not there.

Only if a class is optional for your application (for example, it may come with a plugin), use “\lstinline|Class.forName( <ClassName> )|”\index{java.lang!Class!forName()}; of course you also have to add the proper error handling for the case that the class is not there.

For Reflection\index{reflection}, the situation is similar: you refer to the class element that you want to get access to by its name, provided as a string. The compiler is unable to validate that name, so a potential error will surface earliest at runtime. Refer to the chapter \tqfullvref{sec:Reflection} for more details about Reflection\index{reflection} and its alternatives.

%==============================================================================
% End of file!!
%==============================================================================



%==============================================================================
% End of file!!
%==============================================================================



%==============================================================================
% End of file!!
%==============================================================================

